%% initech-logo.tex -- the Initech logo as original TikZ/PGF drawing commands.
%% Loaded by initech-manual.cls (\input after tikz).  PUB-REF-0001 D-8 as
%% revised by operator directive (WL-0132), Sec 7.1 PK-01 "Initech Cyan",
%% Sec 7.2 PK-02 logo cell.  Bead: initech-ipu5 (identity lane, WL-0132).
%%
%% HONESTY (Law 1): this is an ORIGINAL vector drawing, AUTHORED FROM
%% OBSERVATION of the reference frame's identity material -- the clean logo
%% still and the headquarters-sign still in the private reference library
%% (initech-manuals-ref/initech-identity/, outside this repository), cross-
%% checked against a 1200-pixel third-party rendering of the same mark.
%% Nothing here is traced, extracted or converted from an image: the
%% numbers below are measurements of pixel runs, rounded, and the shapes are
%% drawn from those numbers.  Nothing else from the references is copied.
%%
%% WHAT THE LOGO IS (measured):
%%  * Mark: three identical blocks in a 2:1 dimetric projection (receding
%%    edges rise 1 in 2, i.e. 26.565 deg; vertical edges sqrt(3/2) = 1.2247
%%    of the receding edge's horizontal run -- the true orthographic view
%%    from 30 deg elevation, 45 deg azimuth).  Each block is a cube of edge
%%    2 with a cube of edge 1 removed from its top front corner, so the top
%%    face reads as a chevron and the notch shows a floor and two walls.
%%    Faces: tops and notch floors in Initech Cyan; faces turned to the
%%    viewer's left (and the notch wall parallel to them) near-black; faces
%%    turned to the right (and their notch wall) light grey.  One block at
%%    the back, two in front of it at (-2.5u, -2u) and (+2.5u, -2u) on the
%%    page, where u is the horizontal run of one unit edge: a triangle, one
%%    unit of air between the two front blocks.  Mark: 9u wide x 6.449u tall.
%%  * Logotype: INITECH in capitals of a black-weight, extended grotesque
%%    (Helvetica-class: the C is a round bowl with horizontal terminals, not
%%    the squared Eurostile C), oblique at 0.2065 (11.7 deg), cap height
%%    2.536u, upright set width 14.740u, slanted width 15.264u, tight fitting
%%    (0.25-0.40u between letters).  Stems 0.763u (30 % of the cap height).
%%  * Lockup: mark centred over the slanted logotype's box, 1.471u of air
%%    between cap line and the mark's lowest vertex.  Whole lockup 15.264u
%%    wide x 10.459u high (baseline to apex): aspect 1.4594.
%%  Measurements: clean still 300 x 207 px (u = 19.5 px) and the 1200-px
%%  rendering (u = 62.89 px) agree to 2 %; colours sampled on the cube tops:
%%  RGB 87/177/159 (#57B19F), median of nine samples, flat ink.
%%
%% PUBLIC MACROS (names are stable; other lanes call them):
%%   \InitechLogo[<height>]       colour lockup (covers, spines, boxes, labels)
%%   \InitechLogoBlack[<height>]  one-colour lockup, black + black tints:
%%                                the ONLY form allowed inside a book (D-6)
%%   \InitechMark[<height>]       the three blocks alone, colour
%%   \InitechMarkBlack[<height>]  the three blocks alone, black + tints
%%   \InitechLogotype[<cap ht>]   the letters alone, in the current colour
%%   \InitechLogoAspect           1.4594 (lockup width / height)
%%   \InitechLogoForWidth{<w>}    the lockup height that gives width <w>
%%   Default height: 20 mm (lockup, mark), 5 mm (logotype cap height).
%%   Each macro yields a box whose bottom is the baseline of the logotype
%%   (lockup), the lowest vertex (mark) or the baseline (logotype).
%% Colours used (defined in initech-manual.cls): InitechCyan, InitechBlack,
%%   InitechCubeLight, InitechCubeMono.
\ProvidesFile{initech-logo.tex}[1992/03/01 Initech logo drawing]
\newdimen\im@lgu
\def\InitechLogoAspect{1.459}
%% height for a given width: h = w / 1.4594 = w * 685 / 1000
\newcommand\InitechLogoForWidth[1]{\dimexpr #1 * 685 / 1000\relax}

%% One block, origin = the screen point of the block's back bottom corner
%% (3-D 0,0,0).  Screen (X, Y) = (x - y, z * 1.2247 - (x + y) / 2), unit u.
%% Every face is filled and stroked with a 0.12 pt hairline in its own
%% colour, so abutting faces leave no seam on screen and no halo in print.
\tikzset{im@lgface/.style={fill=#1, draw=#1, line width=0.12pt}}
\newcommand\im@lg@block[1]{\begin{scope}[shift={#1}]
  %% right outer face (x = 2) and the notch wall x = 1: light
  \path[im@lgface=InitechCubeLight] (2,-1) -- (0,-2) -- (0,-0.775) -- (1,-0.275)
    -- (1,0.949) -- (2,1.449) -- cycle;
  \path[im@lgface=InitechCubeLight] (0,0.225) -- (-1,-0.275) -- (-1,0.949)
    -- (0,1.449) -- cycle;
  %% left outer face (y = 2) and the notch wall y = 1: dark
  \path[im@lgface=\im@lg@dark] (-2,-1) -- (0,-2) -- (0,-0.775) -- (-1,-0.275)
    -- (-1,0.949) -- (-2,1.449) -- cycle;
  \path[im@lgface=\im@lg@dark] (0,0.225) -- (1,-0.275) -- (1,0.949) -- (0,1.449)
    -- cycle;
  %% top (z = 2, a chevron) and the notch floor (z = 1)
  \path[im@lgface=\im@lg@top] (0,2.449) -- (2,1.449) -- (1,0.949) -- (0,1.449)
    -- (-1,0.949) -- (-2,1.449) -- cycle;
  \path[im@lgface=\im@lg@top] (0,0.225) -- (1,-0.275) -- (0,-0.775) -- (-1,-0.275)
    -- cycle;
  \end{scope}}
%% The mark, origin = the back block's origin: back block first, then the
%% two front blocks (painter's order).  Extent X -4.5..4.5, Y -4..2.4495.
\newcommand\im@lg@mark{\im@lg@block{(0,0)}\im@lg@block{(-2.5,-2)}%
  \im@lg@block{(2.5,-2)}}
%% The logotype, origin = baseline at the upright left edge of the first I.
%% Upright outlines, then sheared by the measured oblique.
\newcommand\im@lg@letters{\begin{scope}[xslant=0.207]
  \fill (0,0) -- (0.763,0) -- (0.763,2.536) -- (0,2.536) -- cycle;
  \fill (1.161,0) -- (1.908,0) -- (1.908,1.407) -- (2.668,0) -- (3.435,0)
    -- (3.435,2.536) -- (2.687,2.536) -- (2.687,1.145) -- (1.956,2.536)
    -- (1.161,2.536) -- cycle;
  \fill (3.848,0) -- (4.611,0) -- (4.611,2.536) -- (3.848,2.536) -- cycle;
  \fill (5.549,0) -- (6.325,0) -- (6.325,1.884) -- (7.025,1.884)
    -- (7.025,2.536) -- (4.863,2.536) -- (4.863,1.884) -- (5.549,1.884)
    -- cycle;
  \fill (7.314,0) -- (9.413,0) -- (9.413,0.660) -- (8.078,0.660)
    -- (8.078,0.986) -- (9.286,0.986) -- (9.286,1.582) -- (8.078,1.582)
    -- (8.078,1.884) -- (9.398,1.884) -- (9.398,2.536) -- (7.314,2.536)
    -- cycle;
  %% C: outer ellipse centre (10.908,1.272), radii 1.209 x 1.304, cut by
  %% the two horizontal terminals; the counter is four quarter ellipses and
  %% a straight left side.
  \fill (12.072,1.622) arc[start angle=15.563, end angle=347.320,
      x radius=1.209, y radius=1.304]
    -- (11.385,0.986)
    arc[start angle=0, end angle=-90, x radius=0.477, y radius=0.350]
    arc[start angle=270, end angle=180, x radius=0.413, y radius=0.382]
    -- (10.495,1.590)
    arc[start angle=180, end angle=90, x radius=0.413, y radius=0.334]
    arc[start angle=90, end angle=0, x radius=0.477, y radius=0.302]
    -- cycle;
  \fill (12.482,0) -- (13.246,0) -- (13.246,1.018) -- (13.961,1.018)
    -- (13.961,0) -- (14.740,0) -- (14.740,2.536) -- (13.961,2.536)
    -- (13.961,1.654) -- (13.246,1.654) -- (13.246,2.536) -- (12.482,2.536)
    -- cycle;
  \end{scope}}

%% #1 height, #2 lockup height in u (x 10000), #3 drawing, #4 bounding box
\newcommand\im@lg@picture[4]{\begingroup
  \im@lgu=\dimexpr #1 * 10000 / #2\relax
  \begin{tikzpicture}[x=\im@lgu, y=\im@lgu, baseline=(current bounding box.south)]
    \useasboundingbox #4;
    #3
  \end{tikzpicture}\endgroup}
\newcommand\im@lg@colour{\def\im@lg@top{InitechCyan}\def\im@lg@dark{InitechBlack}}
\newcommand\im@lg@mono{\def\im@lg@top{InitechCubeMono}\def\im@lg@dark{InitechBlack}}
\newcommand\im@lg@lockup{%
  \begin{scope}[shift={(7.632,8.010)}]\im@lg@mark\end{scope}%
  \begin{scope}[InitechBlack]\im@lg@letters\end{scope}}
\newcommand\InitechLogo[1][20mm]{{\im@lg@colour
  \im@lg@picture{#1}{104590}{\im@lg@lockup}{(0,0) rectangle (15.264,10.459)}}}
\newcommand\InitechLogoBlack[1][20mm]{{\im@lg@mono
  \im@lg@picture{#1}{104590}{\im@lg@lockup}{(0,0) rectangle (15.264,10.459)}}}
\newcommand\InitechMark[1][20mm]{{\im@lg@colour
  \im@lg@picture{#1}{64495}{\im@lg@mark}{(-4.5,-4) rectangle (4.5,2.449)}}}
\newcommand\InitechMarkBlack[1][20mm]{{\im@lg@mono
  \im@lg@picture{#1}{64495}{\im@lg@mark}{(-4.5,-4) rectangle (4.5,2.449)}}}
\newcommand\InitechLogotype[1][5mm]{%
  \im@lg@picture{#1}{25362}{\im@lg@letters}{(0,0) rectangle (15.264,2.536)}}
\endinput
